% Lösungen Einheit 3 von 5 – Dreiecke nachweisen (zusammenbau v0.1)
\einheitenkopf{Einheit 3 von 5 · Dreiecke nachweisen}

\erg{15}{a) $\overrightarrow{DC}$ \quad b) $|\overrightarrow{CA}| = |\overrightarrow{CB}| = \sqrt{21}$, $|\overrightarrow{AB}| = \sqrt{20}$: gleichschenklig, Schenkel $CA$ und $CB$, Basis $AB$ \quad c) $|\overrightarrow{AB}| = |\overrightarrow{BC}| = |\overrightarrow{CA}| = \sqrt{8}$: gleichseitig. \quad d) $|\overrightarrow{CA}| = |(-2 | -2 | -p)| = \sqrt{8 + p^2}$ und $|\overrightarrow{CB}| = |(2 | 2 | -p)| = \sqrt{8 + p^2}$: gleich für jedes $p$, Basis $AB$. \quad e) Die Dreiecke $ABD_p$ und $ACD_p$ haben bei $A$ einen rechten Winkel (Achsen), gleich lange Katheten $AB$ und $AC$ und die gemeinsame Kathete $AD_p$; sie sind kongruent, also $|BD_p| = |CD_p|$ – Basis ist $BC$. \quad f) $|\overrightarrow{AB}| = |\overrightarrow{AC}| = \sqrt{34}$, Basis $BC$ (nicht Schenkel); $E$ und $F$ haben dieselbe dritte Koordinate $2$, die Grundfläche liegt in der $x_1x_2$-Ebene: $EF$ ist parallel zu ihr. \quad g) $M$ ist der Mittelpunkt von $DE$; $\overrightarrow{FD} \circ \overrightarrow{FE} = 0$, also hat das Dreieck bei $F$ einen rechten Winkel. Nach Thales liegt $F$ auf dem Kreis um $M$ durch $D$ und $E$. \quad h) Symmetrie zu $AB$: $M(3 | m)$; $|MA| = |MC|$: $3^2 + m^2 = (9 - m)^2$ ergibt $m = 4$: $M(3 | 4)$ (nicht der Schwerpunkt $(3 | 3)$) \quad i) $12 + 2\sqrt{36 + (p - 1)^2} = 32$; erst die Wurzel isolieren: $\sqrt{36 + (p - 1)^2} = 10$, $(p - 1)^2 = 64$, $p = 9$ \quad j) $\overrightarrow{OT} = \overrightarrow{PQ} = (-3 | 3 | 2)$, also $T(-3 | 3 | 2)$: $OPQT$ ist ein Parallelogramm, die Dreiecke $OPQ$ und $QTO$ sind kongruent. \quad k) $C$ am Mittelpunkt $M(3 | 3 | 3)$ von $AB$ spiegeln: $D = A + B - C = (1 | 3 | 5)$; $ADBC$ ist dann eine Raute.}
\erg{16}{a) Fehler: Tom vergleicht die Basis mit einem Schenkel. Richtig: $|\overrightarrow{CA}| = |\overrightarrow{CB}| = 3$, also gleichschenklig mit der Basis $AB$. \quad b) Gleichschenklig heißt: zwei Seiten sind gleich lang. Die beiden Seiten, die an der Spitze zusammenstoßen, sind die Schenkel; die dritte ist die Basis und muss nicht gleich lang sein.}
